\documentclass[11pt,a4paper]{article}
% Estilo comum dos PDFs individuais da Sprint 1.
\usepackage[utf8]{inputenc}
\usepackage[T1]{fontenc}
\usepackage[brazilian]{babel}
\usepackage[a4paper,margin=2.35cm]{geometry}
\usepackage{xcolor}
\usepackage{hyperref}
\usepackage{enumitem}
\usepackage{listings}
\usepackage{parskip}
\usepackage{needspace}
\definecolor{azul}{RGB}{20,63,112}
\hypersetup{colorlinks=true,linkcolor=azul,urlcolor=azul}
\setlist[itemize]{leftmargin=*,itemsep=3pt}
\lstset{basicstyle=\ttfamily\small,frame=single,breaklines=true,columns=fullflexible}
\sloppy
\newcommand{\codigo}[1]{\texttt{#1}}
\newcommand{\titulo}[2]{\begin{center}{\LARGE\bfseries #1\par}\vspace{4pt}{\large #2\par}\vspace{8pt}\rule{0.8\textwidth}{0.4pt}\par\vspace{5pt}{\small Matheus Moreli Durão Silva -- Matrícula 12621CBS217\par Cibersegurança - Bacharelado - Noturno\par}\end{center}}
\newcommand{\cardfonte}[1]{\par\smallskip\noindent\textit{Referência: }\href{#1}{Trello - enunciado e checklist do card}.}

\begin{document}\titulo{S1\_02 - Requisitos funcionais}{O que o sistema precisa fazer}\enlargethispage{3\baselineskip}
\section*{Ideia central}
Um requisito funcional (RF) descreve uma ação observável do sistema. A formulação deve conter um verbo e um resultado: ``o sistema deve cadastrar um ativo''. Isso é diferente de escrever apenas ``ativos'', que não informa nenhuma ação. Requisito não funcional (RNF) descreve condição, limite ou qualidade: ``o sistema deve manter os dados em JSON''.
\section*{Requisitos do inventário}
\begin{itemize}
\item RF01: cadastrar ativo com identificador, nome, tipo, responsável e localização.
\item RF02: listar e buscar ativos cadastrados.
\item RF03: atualizar e remover ativo pelo identificador.
\item RF04: registrar vulnerabilidade vinculada a um ativo e informar sua severidade.
\item RF05: visualizar vulnerabilidades e gerar um resumo do inventário.
\item RNF01: validar campos obrigatórios e valores numéricos. RNF02: persistir os dados em JSON. RNF03: operar pelo terminal com mensagens claras.
\end{itemize}
\section*{CRUD e casos de uso}
CRUD é uma sigla para \textit{Create, Read, Update, Delete}: criar, consultar, alterar e excluir. No projeto, cadastro é Create; listar e buscar são Read; atualização é Update; remoção é Delete. O gerenciamento de vulnerabilidades complementa o CRUD porque registra um risco relacionado a um ativo.

Exemplo de caso de uso: \textbf{Cadastrar ativo}. Ator: usuário do inventário. Fluxo: escolhe cadastro, informa os campos, o sistema valida os valores, guarda o ativo e confirma o resultado. Resultado esperado: o ativo aparece em uma consulta posterior.
\par\medskip
\textbf{Buscar ativo}. Ator: analista. Pré-condição: inventário carregado. Fluxo: escolhe busca por ID ou nome, informa o critério e recebe os ativos correspondentes. Alternativa: identificador inexistente produz ``Ativo não encontrado'', sem criar dados. Pós-condição: coleção preservada.
\par\medskip
\textbf{Remover ativo}. Ator: responsável pelo inventário. Pré-condição: ativo existente. Fluxo: informa o ID, confere o alvo e confirma. O sistema exclui o ativo e suas vulnerabilidades. Alternativa: cancelamento mantém os dados. Pós-condição: o ID excluído não aparece nas consultas.
\section*{Como isso aparece no código}
No \codigo{app.py}, as opções do menu são a tradução dos RFs. Funções separadas como \codigo{cadastrar\_ativo} e \codigo{buscar\_ativo} tornam possível relacionar uma função do código a uma necessidade do sistema. Se um requisito não estiver associado a um fluxo, ainda falta implementação ou especificação.
\section*{Checklist demonstrado}
Foram estudados: conceito de RF; diferença entre RF e RNF; identificação de funcionalidades; caso de uso; critérios de aceitação; relação entre requisitos e CRUD. Cada item acima pode ser revisado perguntando: ``qual ação o usuário realiza, qual dado informa e qual resultado consegue observar?''
\cardfonte{https://trello.com/c/nPnrvQad}
\end{document}
